% Intended LaTeX compiler: pdflatex
\documentclass[a4paper,12pt]{article}
\usepackage[utf8]{inputenc}
\usepackage[T1]{fontenc}
\usepackage{graphicx}
\usepackage{longtable}
\usepackage{wrapfig}
\usepackage{rotating}
\usepackage[normalem]{ulem}
\usepackage{amsmath}
\usepackage{amssymb}
\usepackage{capt-of}
\usepackage{hyperref}
\usepackage{\string~/.emacs.d/common}
\author{mahmood}
\date{\today}
\title{algorithms homework 6}
\hypersetup{
 pdfauthor={mahmood},
 pdftitle={algorithms homework 6},
 pdfkeywords={},
 pdfsubject={},
 pdfcreator={Emacs 28.2 (Org mode 9.5.5)}, 
 pdflang={English}}
\begin{document}

\maketitle
homework on \href{20230102173117-dynamic_programming.org}{dynamic programming}\\
\begin{note}
i kept coming back to the problems trying to correct/improve the solutions many times, but at some point i just gave up and handed the homework over, would love feedback though!\\
\end{note}
\begin{question}
provide an algorithm that is as efficient as possible that solves the following problem:\\
\uline{input}: an array of \href{calculus2.org}{natural} numbers\\
\uline{output}: a list of k pairs (\(k \le n/2\)) of indicies of the form i,i+1 such that none of the indicies of the array appears more than once and the sum \(\sum A[i]-A[i+1]\) is maximal\\
example: for the input 12,11,2,4,1,10,8,12 the output should be (11,2),(4,1),(10,8)\\
write an algorithm to find the pairs of indicies as well as the sum\\
\begin{subquestion}
describe the algorithm in words\\
\begin{answer}
its all in the next subquestion\\
\end{answer}
\end{subquestion}
\begin{subquestion}
write pseudocode for the algorithm\\
\begin{answer}
first we propose a solution, then optimize it using dynamic programming\\
\includesvg{/home/mahmooz/brain/out/2CcvjGm}\\

to make this simpler for implementation and for further optimization we rewrite the algorithm and use an additional argument, \(i\), thats an index whose initial value should be \(i=|A|-1\), because thats how the algorithm i've written works\\
\includesvg{/home/mahmooz/brain/out/oPeEeuE}\\

this algorithm has exponential time complexity, we could precompute the values that would be returned by \(\textsc{Question1-Alg-Recurse}\) into an array B, such that \(B[i] := \textsc{Question1-Alg-Recurse}(A, i)\), we use another array C to keep track of the pairs we pick, only because we are required to do so\\
\includesvg{/home/mahmooz/brain/out/N4Ymsmc}\\

by running \(\textsc{Question1-Alg-Efficient}\), \(C\) would contain the pairs of indicies and \(B[|A|-1]\) the sum we're looking for, this algorithm runs in \(\Theta(n)\)\\
\begin{note}
there is something i couldnt get quite right with \(C\)\\
i think this is wrong, there is something wrong with the way im keeping record of the pairs in \(C\) but im too lazy/exhausted to figure it out\\
\end{note}
\end{answer}
\end{subquestion}
\begin{subquestion}
show the \href{20221104220603-algorithm_correctness.org}{correctness} of the algorithm you proposed\\
\begin{answer}
we prove the correctness of the algorithm using \href{20220707193301-mathematical_induction.org}{induction} on the size of the input array\\
\begin{enumerate}
\item base step: for \(|A|=2\), the solution, \(B[1]\), would equal \(A[1]-A[2]\), and C to a singleton of \((1,2)\) which is the only pair in the array so it is the solution\\
\item the \href{20230108230724-induction_hypothesis.org}{induction hypothesis}: we assume truth for \(|A|=n\)\\
\item the \href{20230108230748-induction_step.org}{induction step}, we show truth for \(|A|=n+1\)\\
on the last iteration of the for loop \(4,\dots,n+1-1\), we would be checking which of the values \(A[i-1]-A[i]+B[i-3]\),\(A[i]-A[i+1]+B[i-2]\) is bigger, the first corresponding to a \href{20221204105120-combinatorics.org}{combination} of pairs whose sum would've been maximal up until \(A[i-2]\), and the second corresponding to a combination of pairs whose sum would've been maximal up until \(A[i-1]\), of those two, we pick one and add to it yet another pair only if the addition of that pair makes the sum stay or become bigger than the sum of the other combination with the other new pair added to it, so the last sum in B would still be maximal, and the pairs in C, likewise, would be corresponding to the newly found maximal sum\\
\end{enumerate}
\begin{note}
im not sure of this proof, should reevaluate it\\
\end{note}
\end{answer}
\end{subquestion}
\begin{subquestion}
what is the \href{20221130014441-time_complexity.org}{time complexity} of the algorithm?\\
\begin{answer}
its simply a single for loop with no added complexity whatsoever, so its just \(\Theta(n)\)\\
\end{answer}
\end{subquestion}
\label{org6bb891d}
\end{question}
\begin{question}
given an array \(A\) of integers of size \(n \times n\), a "legal" route in this array is one that starts at \(A[1,1]\) and ends at \(A[n,n]\), a legal move is one from \(A[i,j]\) to \(A[i+1,j]\) or \(A[i,j+1]\), the cost of a legal route is the sum of the elements it contains such that the maximal element is added twice\\
write an algorithm that finds the maximal cost of a legal route in a given array\\
\begin{subquestion}
write the algorithm\\
\begin{answer}
we use an array B, such that \(B[x,y]\) would hold the maximal sum of all legal routes from \(A[1,1]\) to \(A[x,y]\)\\
we use another array C, such that \(C[x,y]\) would hold the maximal element from \(A[1,1]\) to \(A[x,y]\)\\
\includesvg{/home/mahmooz/brain/out/EnsRPuP}\\
\end{answer}
\end{subquestion}
\begin{subquestion}
show the correctness of the algorithm\\
\begin{answer}
we propose the following \href{20221104221055-loop_invariant.org}{loop invariant} for the inner loop:\\
\begin{lemma}
at the end of each \(j\)'th iteration, \(C[i,j]\) (\(i\) is constant here and denotes the \(i\)'th row) contains the maximal number in the optimal route \(A[1,1]\) from to \(A[i,j]\) under the constraint that \(A[i,j]\) is in that route, and \(B[i,j]\) contains the sum of the numbers in that route with the maximal number added twice\\
\begin{proof}
\begin{note}
this proof is incorrect, in the base case im assuming \(B[i-1,1]\) is the correct sum of the route \(A[1,1]\) to \(A[i-1,1]\), im not sure how to correct this :/\\
\end{note}
\begin{enumerate}
\item base step: for \(j=1\), we set \(B[i,1]\) to \(A[i,1]+B[i-1,1]\) as it should be equal to the sum of the optimal route from \(A[1,1]\) to \(A[i,1]\), and we only subtract the previously maximal number from the sum \(B[i,1]\) and add \(A[i,1]\) again if \(A[i,1]\) is bigger, and \(C[i,1]\) would also be updated accordingly\\
\item the \href{20230108230724-induction_hypothesis.org}{induction hypothesis}: we assume truth for \(j=k\)\\
\item the \href{20230108230748-induction_step.org}{induction step}, we show truth for \(j=k+1\)\\
we set \(B[i,k+1]\) to \(A[i,k+1]+B[i,k]\) as it should be equal to the sum of the optimal route from \(A[1,1]\) to \(A[i,k+1]\), and we only subtract the previously maximal number from the sum \(B[i,k+1]\) and add \(A[i,k+1]\) again if \(A[i,k+1]\) is bigger, and \(C[i,k+1]\) would also be updated accordingly\\
\end{enumerate}
\end{proof}
\end{lemma}
we propose the following loop invariant for the outer loop:\\
\begin{lemma}
after each \(i\)'th iteration, \(B[i,j]\) would equal the required sum (with the maximal number added twice) of the optimal legal route from \(A[1,1]\) to \(A[i,j]\) and \(b_{\max}\) is equal to the maximum sum in \(B\)\\
\end{lemma}
\end{answer}
\end{subquestion}
\begin{subquestion}
find the time complexity of the algorithm\\
\begin{answer}
the algorithm executes two for loops with constant execution time on each iteration, so in total \(\Theta(n^2)\), which is the \href{20221130014441-time_complexity.org}{best conceivable runtime}, because the algorithm has to touch each element in the array atleast once for it to decide which route is the optimal one (i.e. the one with the highest sum)\\
\end{answer}
\end{subquestion}
\end{question}
\begin{question}
given an array of numbers of size \(n \times n\), find the maximal sum of any \(n\) elements in the array, under the constraints that from each row only one element is picked, and each element has to be picked in the same column as the one before, or the column after or the column before\\
\begin{subquestion}
explain the algorithm in words\\
\begin{answer}
we use an array \(B\) such that \(B[x,y]\) is the solution to the problem in the space \(A[1,\dots,n,\ 1,\dots,y]\) under the constraint that \(A[x,y]\) is part of the solution (so that the solutions that dont include \(A[x,y]\) arent taken into account)\\
then we iterate through the data we found to find the combination of elements with the biggest sum\\
\end{answer}
\end{subquestion}
\begin{subquestion}
write pseudocode that solves the problem\\
\begin{answer}
\includesvg{/home/mahmooz/brain/out/anL9Kp8}\\
\end{answer}
\end{subquestion}
\begin{subquestion}
show the correctness of the algorithm\\
\end{subquestion}
\begin{subquestion}
find the time complexity of the algorithm\\
\begin{answer}
\(\Theta(n^2)\)\\
\end{answer}
\end{subquestion}
\end{question}
\begin{question}
given an array of numbers of size \(n\), find the maximal length of a subarray of sorted elements in ascending order under the constraint that the difference between the indicies of the subarray is at most 2\\
\begin{answer}
we use an additional array \(B\), such that \(B[i]\) equals the length of the optimal subarray in the space \(A[1,\dots,i]\) under the constraint that \(B[i]\) is in the subarray, the solution would be \(\max B\)\\
\includesvg{/home/mahmooz/brain/out/p19PEiW}\\

this algorithm runs in \(\Theta(n)\)\\
\end{answer}
\end{question}
\begin{question}
asked to prove a \href{./20230102173117-dynamic\_programming.org}{given lemma}\\
\begin{answer}
we assume in contradiction that neither \(A[n]\) nor \(A[n-1]\) are in the optimal solution \(B\)\\
since the numbers are positive, we know for sure that \(\sum B + A[n] > \sum B\), and since \(\sum B + A[n]\) is a legal solution to the problem (because the closest element that could be in the array \(B\) to \(A[n]\) is \(A[n-2]\) which isnt next to \(A[n]\)), we've arrived at a contradiction as we assumed \(B\) was the optimal solution\\
and so one of \(A[n]\) and \(A[n-1]\) have to be in the optimal solution\\
\end{answer}
\end{question}
\end{document}